\chapter{GemCrush : ce qui change quand un jeu grossit}

Les trois jeux de plateau tenaient en quatre fichiers. GemCrush en a vingt-deux
et huit mille lignes --- pour un jeu d'alignement dont la règle tient en une
phrase : \emph{échangez deux gemmes voisines, alignez-en trois ou plus, elles
disparaissent}.

Ce chapitre explique où passent les huit mille lignes. La réponse n'est pas « la
règle est compliquée » : elle est que tout ce qui entoure la règle --- écrans,
progression, animation, son, didacticiel --- coûte bien plus que la règle
elle-même. C'est la leçon la plus utile du chapitre.

\section{Le découpage}

\begin{center}
\begin{tabular}{@{}lll@{}}
\toprule
\textbf{Fichier} & \textbf{Rôle} & \textbf{Famille} \\
\midrule
\texttt{NkGem} & la gemme : couleur, forme, état & règle \\
\texttt{NkGemBoard} & le plateau : chutes, échanges & règle \\
\texttt{NkMatchFinder} & trouver les alignements & règle \\
\texttt{NkGemLevels} & les trente niveaux & données \\
\texttt{Ui/NkGemArt} & dessiner une gemme & écran \\
\texttt{Ui/NkGemScreens} & menu, carte, partie & écran \\
\texttt{Ui/NkGemHud} & score, objectifs, coups restants & écran \\
\texttt{Ui/NkGemSplash} & l'ouverture & écran \\
\texttt{Ui/NkGemTutorial} & le didacticiel & écran \\
\texttt{Ui/NkGemTheme} & les couleurs & écran \\
\texttt{Ui/NkGemAudio} & les sons, synthétisés & son \\
\bottomrule
\end{tabular}
\end{center}

\begin{nkretenir}
\textbf{Le même principe qu'aux jeux de plateau, à une échelle plus grande} : les
trois premiers fichiers sont la règle et ne dessinent rien ; tout ce qui est sous
\texttt{Ui/} dessine et ne décide de rien.

Ce qui a changé, ce n'est pas le principe : c'est le \emph{nombre} de fichiers de
chaque côté. Quand un jeu grossit, c'est presque toujours le côté écran qui
grossit.
\end{nkretenir}

\section{Les quatre écrans}

\begin{nkcode}{Applications/Gemcrush/src/Gemcrush/Ui/NkGemScreens.h}
enum class GemScreen : uint8 { Splash, Menu, Map, Game };
\end{nkcode}

Une énumération, et une seule variable qui dit où l'on est. Le \texttt{OnRender}
principal aiguille dessus.

\begin{nkretenir}
\textbf{Un jeu qui a plus de deux écrans a besoin d'une machine à états
explicite.} Sans elle, on finit avec six booléens --- \texttt{menuOuvert},
\texttt{enPartie}, \texttt{surLaCarte}\dots --- qui peuvent être vrais en même
temps, ce qui n'a aucun sens et arrive tout de même.

Une énumération rend l'état impossible à contredire : on est dans un écran, et un
seul.
\end{nkretenir}

\section{Le son, sans un seul fichier}

\texttt{NkGemAudio} \textbf{synthétise} ses sons : il calcule les échantillons au
lieu de lire un \texttt{.wav}.

\begin{nkretenir}
\textbf{Aucun asset, et c'est un choix documenté du jeu} : « GemCrush ne charge
aucun fichier à l'exécution [...] à l'identique sur les six plateformes sans
chaîne de copie d'assets ».

Le prix : les sons sont simples. Le gain : rien à copier, rien à empaqueter, rien
à précharger sur le Web --- et un jeu qui démarre pareil partout. Pour un jeu de
puzzle, l'échange est bon ; pour un jeu à musique, il ne le serait pas.
\end{nkretenir}

\section{Le didacticiel, et ce qu'il révèle}

\texttt{NkGemTutorial} est un fichier entier pour apprendre au joueur à échanger
deux gemmes.

\begin{nkretenir}
\textbf{Ce n'est pas du superflu.} Un jeu dont la règle tient en une phrase n'est
pas un jeu qu'on comprend en une phrase : il faut montrer le geste, une fois, au
bon moment, puis se taire.

Le didacticiel de GemCrush s'efface seul et devient muet au niveau 7 --- il ne
répète pas ce qui est acquis.
\end{nkretenir}

\section{Ce que son banc mesure}

\begin{center}
\begin{tabular}{@{}ll@{}}
\toprule
\textbf{Cas} & \textbf{Ce qu'il prouve} \\
\midrule
carte : dessin et clic d'accord & le clic tombe là où le n\oe ud est dessiné \\
les 30 niveaux sont déterministes & même graine, même plateau \\
étoiles : accord affichage/verdict & ce qui s'affiche est ce qui est calculé \\
ouverture : 4,1 s, sautable en 2 appuis & le splash ne bloque pas \\
didacticiel : s'efface seul, muet au niveau 7 & il ne harcèle pas \\
progression : écrite, relue à neuf, effaçable & la sauvegarde tient \\
\bottomrule
\end{tabular}
\end{center}

\begin{nkretenir}
\textbf{Deux de ces cas ne parlent pas de la règle du jeu.} « le clic tombe là où
le n\oe ud est dessiné » et « ce qui s'affiche est ce qui est calculé » vérifient
un \emph{accord entre deux chemins}.

C'est la forme de banc la plus rentable quand un jeu grossit : dès que deux
chemins doivent produire le même verdict --- le dessin et le clic, l'affichage et
le calcul --- l'un est le contrôle de l'autre. S'ils appliquent chacun leur copie
de la règle, rien ne signale le jour où elles divergent.
\end{nkretenir}

\begin{nkpiege}
\textbf{« La progression est relue à neuf » n'est pas la même chose que « la
progression se relit ».} Un aller-retour qui compare le fichier à lui-même
certifie une amputation reproductible : si l'enregistrement perd un champ, la
relecture le perd aussi, et le contrôle reste vert.

Un contrôle d'aller-retour porte \emph{deux} exigences : la stabilité (le second
enregistrement égale le premier) et la \textbf{conservation} (le contenu attendu
est encore là, champ par champ). C'est la seconde qu'on oublie.
\end{nkpiege}

\begin{nkbilan}
Vingt-deux fichiers pour une règle d'une phrase : ce sont les écrans, la
progression, le didacticiel et le son qui pèsent, pas la règle. Le découpage
reste celui des jeux de plateau --- la règle ne dessine pas, l'écran ne décide
pas --- mais le côté écran grossit. Au-delà de deux écrans, une énumération
explicite remplace les booléens qui se contredisent. Et les bancs les plus utiles
d'un gros jeu vérifient l'\emph{accord entre deux chemins}, pas la règle.
\end{nkbilan}

\begin{nkquestions}
\begin{enumerate}
\item Quels fichiers de GemCrush sont du côté « règle » ? Comment le sait-on ?
\item Pourquoi une énumération d'écrans plutôt que plusieurs booléens ?
\item Pourquoi GemCrush synthétise-t-il ses sons au lieu de lire des fichiers ?
      Quel est le prix de ce choix ?
\item Qu'est-ce qu'un banc « d'accord entre deux chemins » ? Donnez-en deux.
\item Quelles sont les deux exigences d'un contrôle d'aller-retour ?
\end{enumerate}
\end{nkquestions}

\begin{nkpratique}
Écrivez le c\oe ur d'un jeu d'alignement : un plateau $8\times8$ de couleurs, une
fonction qui trouve tous les alignements de trois ou plus, et une fonction qui
fait chuter les gemmes pour combler les trous.

\textbf{Sans dessiner quoi que ce soit.} Vérifiez par un banc : un plateau
fabriqué à la main dont vous connaissez les alignements attendus, et le plateau
d'après la chute, comparé case par case.
\end{nkpratique}

\begin{nkdemo}
Prouvez que votre détection d'alignements ne trouve pas ce qui n'existe pas.

Construisez un plateau \emph{sans aucun alignement} et vérifiez que la fonction
rend zéro. Puis --- et c'est la partie qui compte --- construisez un plateau avec
deux gemmes identiques côte à côte seulement, et vérifiez qu'elle rend encore
zéro.

\textbf{Pourquoi ce second cas} : une fonction qui compterait les paires au lieu
des triplets passerait le premier test et raterait le second. Un cas qui doit
rendre zéro prouve autant qu'un cas qui doit trouver.
\end{nkdemo}

\begin{nklogique}
GemCrush a 30 niveaux « déterministes : même graine, même plateau ».

Expliquez pourquoi c'est indispensable à un banc. Puis répondez à ceci : si le
jeu tirait ses gemmes avec une graine prise sur l'horloge, quels cas du tableau
ci-dessus deviendraient impossibles à écrire ?

Enfin : un jeu déterministe est-il un jeu prévisible pour le joueur ? Justifiez
la différence.
\end{nklogique}
